
@@M0@@

D'après 1.3,
pour @@M1@@ assez grand, les conditions suivantes seront remplies :
@@M2@@
@@M3@@

Il en résulte que pour @@M4@@ assez grand, ces éléments sont tous dans @@M5@@, et vérifient la relation écrite plus haut, et que les suites~:
@@M6@@
sont exactes. Cela entraine que dans @@M7@@, @@M8@@ et @@M9@@ ont même classe.

1.11. Supposons donnée une équation de dépendance linéaire sur @@M10@@ entre @@M11@@ :
@@M12@@

On peut écrire chacun des @@M13@@ comme différence des classes de deux faisceaux, soit : @@M14@@, @@M15@@ et @@M16@@ étant localement libres de type fini sur @@M17@@. Posons
@@M18@@

L'égalité (1.11.1) s'écrit alors :
@@M19@@

D'après le lemme 1.10, on obtient pour @@M20@@ assez grand :
@@M21@@
